\section{Theory: the demotion gap and a two-term account}
\label{sec:theory}

\subsection{Estimand}
\label{sec:estimand}

\paragraph{Objective and noising.}
We first define the training objective whose gradients are under test.
A flow-matching model is trained to predict the constant-velocity
transport between a clean latent and pure noise. With clean latent $x$,
noise $\epsilon \sim \mathcal{N}(0, I)$, and a noise level
$\sigma \in (0,1]$, the noised input is the linear interpolation
$z_\sigma = (1-\sigma)\,x + \sigma\,\epsilon$ (at $\sigma \to 0$ the
input is the clean latent, at $\sigma = 1$ it is pure noise), and the
regression target is the velocity $v = \epsilon - x$ that carries
$z_\sigma$ along the interpolation. With caption condition $c$, the loss
is $\|\hat v_\theta(z_\sigma, \sigma, c) - v\|^2$, and gradients are
taken with respect to the LoRA adapter parameters only; the backbone is
frozen.

\paragraph{Demotion.}
A \emph{demotion} $e_0 {\to} e$ replaces the native latent
with a coarser-grid one, produced exactly as a scale-wise trainer would
produce it; $e_0$ denotes the route's source (native) edge
and $e$ its target edge. With $y$ the image at its native (tier-$e_0$)
resolution, $\mathcal{E}$ the VAE encoder, and $\mathcal{R}_e$ the
pixel-space downscale to the tier-$e$ bucket (native aspect ratio
preserved),
\begin{equation*}
x_{\src} \;=\; \mathcal{E}(y)
\quad\longrightarrow\quad
x_{\dem} \;=\; \mathcal{E}\big(\mathcal{R}_{e}(y)\big),
\end{equation*}
a latent with proportionally fewer spatial tokens; we do not downsample
in latent space. Everything else about the training step is held fixed:
the caption $c$, the $\sigma$ under test, and the noise draw structure
(the same stratified draw scheme, with $\epsilon$ drawn at the tier-$e$
shape). The demoted step then trains on
$z_\sigma = (1-\sigma)\,x_{\dem} + \sigma\,\epsilon$ with target
$v = \epsilon - x_{\dem}$, so the substitution changes both what the network
sees and what it is asked to predict, but only through the spatial
grid; nothing else in the step references the resolution.

\paragraph{The estimand.}
Three \emph{arms} recur throughout: $\src$, the source arm,
which trains on the native-resolution latent $x_{\src}$; $\dem$,
the demoted arm, which trains on $x_{\dem}$; and $\reenc$,
a control arm that decodes and re-encodes the native latent at native
resolution: it pays the VAE round trip that demotion necessarily pays,
with no grid change. For an arm
$a \in \{\src, \reenc, \dem\}$, let
$\bar g_a(\sigma) = \mathbb{E}_\epsilon[\nabla_\theta \mathcal{L}]$ be
the population per-bin mean adapter gradient of a query (image,
caption). The \emph{population demotion distance} of the route
$e_0 {\to} e$ is
\begin{equation}
d_{e_0 \to e}(\sigma) \;=\; 1 - \cos\!\big(\bar g_{\src}(\sigma),\,
\bar g_{\dem}(\sigma)\big)
\;=\; \tfrac{1}{2}\,\big\|\hat g_{\src} - \hat g_{\dem}\big\|^2,
\label{eq:gapdef}
\end{equation}
with $\hat g$ the unit-normalized gradients;
$d_{e_0 \to e} \ge 0$ always. One
notational convention holds from here on: arm labels ($\src$, $\reenc$,
$\dem$) mark \emph{whose} gradient, latent, or Jacobian an object is;
the subscript on every route-level scalar is the \emph{route},
written in full, $d_{e_0 \to e}$, where several sources appear, and
abbreviated to the target edge alone --- an edge value, never an arm
label, so it instantiates numerically ($d_e$, $\Floor_{1120}$,
$\mathrm{Resid}_{768}$) --- wherever the source edge is fixed by
context, as it is below. The
control's distance $d_{\reenc}$ --- Eq.~\ref{eq:gapdef} with
$\bar g_{\reenc}$ in place of $\bar g_{\dem}$ --- carries its arm label
instead: the control changes no grid and has no target edge. The
\emph{reported} quantity of every bin-resolved map and every fit in this
paper is the \emph{re-encoding excess}
$\gap_e(\sigma) := d_e(\sigma) - d_{\reenc}(\sigma)$, the demotion
distance beyond what the control already pays; $\gap_e$ can be
negative, $d_e$ cannot. Endpoint and floor tables report debiased
estimates of $d_e$ with the control as its own row; each table states
which object it reports.

In a batched training scenario the same definition admits two
aggregations over a probe set, and the aggregation operator is part of
the estimand: per-example (cosine per image, then mean) or
batch-aggregate (gradients summed over the batch before the cosine;
the object SGD actually follows). The two are related by a
decomposition, not an ordering. Writing each image's demotion-induced
gradient disagreement as a coherent mean $b$ plus a zero-mean
idiosyncratic deviation with covariance $\Sigma_\eta$, the
batch-aggregate distance at batch size $B$ is, to second order, an
intercept plus a $1/B$ term,
\begin{equation*}
\mathbb{E}[d_B] \;\approx\;
\frac{\|P^{\perp} b\|^2}{2\|\mu\|^2}
\;+\;
\frac{\mathrm{tr}\!\big(P^{\perp} \Sigma_\eta P^{\perp}\big)}{2B\|\mu\|^2},
\end{equation*}
with $\mu$ the mean source gradient and $P^{\perp}$ the projector off
$\hat\mu$: the coherent share never averages out. The per-example map
is therefore \emph{not} a lower bound for larger batches --- monotone
improvement with $B$ would require an iid zero-mean disagreement
model we have not verified. These are different
estimands, and no coefficient fitted on one is guaranteed to transfer
to the other; the application section publishes one map for each
(\S\ref{sec:map}).

Finally, any estimate of $d_e(\sigma)$ compares finite-draw estimates
of population directions, so every claim carries the instrument's
resolution. For a route whose true excess is zero, the paired estimate
$\bgap$ fluctuates with standard error
$\mathrm{SE}(N, D, \sigma\text{-bin})$; its one-sided 95\% upper bound
defines the \emph{certification resolution} $\epsstar$ (the
smallest margin certifiable even for a truly safe route) and ``safe
at margin $\varepsilon$'' as the non-inferiority statement
\begin{equation}
\text{UB}_{95}\!\left(\bgap\right) \;<\; \varepsilon,
\qquad
\varepsilon \;\ge\; \epsstar(N, D, \text{bin})
\;=\; 1.645\,\mathrm{SE}(N, D, \text{bin}),
\label{eq:safe}
\end{equation}
not a hard bound. Throughout, ``no certifiable window'' means that a
route's upper bounds do not cross below $\varepsilon$ at the current
$(N, D)$, while ``certifiably unsafe'' means that its lower bound stays
above it.

\subsection{Data and graph perturbations}
\label{sec:branches}

We now derive the perturbation demotion applies to the gradient.
Differentiating the loss of \S\ref{sec:estimand} in the adapter
parameters gives, per draw and up to a constant,
\begin{equation*}
g \;=\; \nabla_\theta\, \tfrac{1}{2}
\big\|\hat v_\theta - v\big\|^2 \;=\; J^{\!\top} r,
\end{equation*}
with $r = \hat v_\theta - v$ the prediction residual and
$J = \partial \hat v_\theta / \partial \theta$ the Jacobian through
which the compute graph turns a residual into an adapter gradient.
Demotion changes exactly the two factors of this product, and nothing
else about the map from (image, noise, caption) to adapter gradient:
the residual, because it changes the data the objective consumes, and
the Jacobian, because it changes the compute graph the gradient flows
through. Averaging over draws and expanding the product of
the perturbed factors (an elementary but bookkeeping-heavy step,
written out in Appendix~\ref{app:branches} together with what exactly
the remainder collects), the perturbation
$\delta g = \bar g_{\dem} - \bar g_{\src}$ decomposes as
\begin{equation}
\delta g \;=\;
\underbrace{J^{\!\top}\Delta \bar r(\sigma)}_{\text{data branch}}
\;+\; \underbrace{\Delta J^{\!\top}\, \bar r}_{\text{graph branch}}
\;+\; \underbrace{\Delta J^{\!\top}\Delta \bar r + \cdots}_{\text{remainder }R_e},
\label{eq:branches}
\end{equation}
where:

\begin{enumerate}
\item \textbf{the data branch $J^{\!\top}\Delta\bar r$}: demotion changes the \emph{data}
the objective consumes; spatial-frequency bands above the coarse
grid's Nyquist limit are absent from the input \emph{and} from the
target $v = \epsilon - x$, and the re-encode perturbs what survives,
yielding a perturbation $\Delta \bar r(\sigma)$ of the mean prediction
residual.
Its input-mediated part shrinks as the model's reliance on input detail
shrinks with $\sigma$; its target-mediated part does not: perturbing the
image by $\delta x$ perturbs the residual by
$D_z f\,[(1{-}\sigma)\,\delta x] + \delta x$, and the second term
carries unit weight at every $\sigma$. At $\sigma{=}1$ the input is pure
noise on every grid, so the input-mediated part vanishes in
distribution; the endpoint is nevertheless \emph{not} data-free by
construction, and whether the surviving target-mediated share is
resolvable is an empirical question, discharged in \S\ref{sec:floor}
(it is real and large, but demotion perturbs it almost entirely
\emph{parallel} to the native gradient, so the angular estimand does
not see it).
\item \textbf{the graph branch $\Delta J^{\!\top}\bar r$}: token count, rotary phase
density, and sequence-length-dependent normalization, collected as an
operator perturbation $\Delta J$ that is $\sigma$-\emph{independent},
because the grid change does not reference $\sigma$.
\end{enumerate}

One qualification is worth stating precisely: $J_{\src}$ and
$J_{\dem}$ act on different grids, so the two branches are not literal
matrix products across arms; they are defined \emph{operationally}, as
the two independent perturbation directions in the shared
adapter-parameter space where both gradients live, isolated by the
designated probes of \S\ref{sec:floor}.

This decomposition already shows why input-level spectral sufficiency
cannot bound the gradient. Input noise masking alone provides no bound
on expected-gradient disagreement, because systematic target and graph
terms survive noise averaging: the quantity training consumes is
$\mathbb{E}_\epsilon[\nabla_\theta \mathcal{L}]$, the target carries the
clean image at coefficient one at every $\sigma$, and $\Delta J$ does
not reference the spectrum at all. Input-level sufficiency licenses
inference-time substitution; it does not extend to training-gradient
equivalence.

\subsection{Angular geometry: the exact form and its expansions}
\label{sec:geometry}

It remains to read the perturbation of Eq.~\ref{eq:branches} in the
units of the estimand, Eq.~\ref{eq:gapdef}: a pure exercise in
cosine geometry, whose forms below are what every measurement in
\S\ref{sec:evidence} is fit against (step-by-step derivation in
Appendix~\ref{app:geometry}). Writing
$u^{\perp} = u - (\hat g_{\src}^{\!\top}u)\,\hat g_{\src}$ for the
component of a vector orthogonal to $\bar g_{\src}(\sigma)$, split the
perturbation into a rescaling and a rotation:
$\kappa_\parallel = \hat g_{\src}^{\!\top}\delta g \,/\, \|\bar g_{\src}\|$
and $\kappa_\perp = \|\delta g^{\perp}\| \,/\, \|\bar g_{\src}\|$. Substituting
$\bar g_{\dem} = \bar g_{\src} + \delta g$ into
Eq.~\ref{eq:gapdef} then gives, \emph{exactly},
\begin{equation}
d_e \;=\; 1 - \frac{1}{\sqrt{1 + \kappa_{\mathrm{eff}}^2}}\,,
\qquad
\kappa_{\mathrm{eff}} \;=\; \frac{\kappa_\perp}{1 + \kappa_\parallel}\,:
\label{eq:exact}
\end{equation}
the gap is a saturating function of a single rotation-to-scale ratio:
only the orthogonal component rotates the gradient direction, and
a parallel component merely renormalizes it. Taylor-expanding
Eq.~\ref{eq:exact} at small $\kappa$
($d_e = \tfrac{1}{2}\|\delta g^{\perp}\|^2/\|\bar g_{\src}\|^2 + \cdots$) and
substituting the branch decomposition of Eq.~\ref{eq:branches} into
the square
gives the \emph{four-term angular expansion}, which, like
Eq.~\ref{eq:exact}, is derived unconditionally:
\begin{equation}
d_e(\sigma) \;\approx\;
\underbrace{\frac{\big\|(J^{\!\top}\Delta\bar r)^{\perp}\big\|^2}{2\,\|\bar g_{\src}\|^2}}_{\text{data share }S_e}
\;+\;
\underbrace{\frac{\big\|(\Delta J^{\!\top}\bar r)^{\perp}\big\|^2}{2\,\|\bar g_{\src}\|^2}}_{\text{graph share }F_e}
\;+\;
\underbrace{\frac{\big\langle (J^{\!\top}\Delta\bar r)^{\perp},\, (\Delta J^{\!\top}\bar r)^{\perp}\big\rangle}{\|\bar g_{\src}\|^2}}_{\text{projected interaction }I_e}
\;+\; R'_e,
\label{eq:fourterm}
\end{equation}
where $R'_e$ now has an exact characterization
(Appendix~\ref{app:geometry}): it collects the remainder branch $R_e$,
the beyond-quadratic terms of Eq.~\ref{eq:exact}, and the
parallel-component correction. $S_e$ and $F_e$ are non-negative; $I_e$
can carry either sign, a fact that matters in
\S\ref{sec:confronted}. The expansion's validity domain deserves
stating where it is introduced: at the harshest measured operating
point ($d = 0.30$ at $512$, $\sigma{=}1$) the implied perturbation is
$\kappa_{\mathrm{eff}} = \sqrt{(1-d)^{-2} - 1} \approx 1.02$ --- not
small --- so quadratic-order reads on that route are indicative, and
the exact form (Eq.~\ref{eq:exact}) carries the quantitative claims.

\paragraph{From the expansion to the excess.}
The reported object is the excess over the control
(\S\ref{sec:estimand}): let
$\gap_e(\sigma) = d_e(\sigma) - d_{\reenc}(\sigma)$. The control's own
expansion is degenerate: the $\reenc$ arm changes no grid, so
$\Delta J_{\reenc} = 0$, and both control terms built from it vanish
identically --- its graph share $F_{\reenc}$ and its projected
interaction $I_{\reenc}$ each carry a factor
$(\Delta J_{\reenc}^{\!\top}\bar r)^{\perp} = 0$ --- leaving only the
re-encode perturbation in its data branch. Writing
Eq.~\ref{eq:fourterm} for each arm and subtracting term by term,
\begin{align*}
\gap_e(\sigma)
\;&=\; d_e(\sigma) \;-\; d_{\reenc}(\sigma)\\
&\approx\;
\big(S_e - S_{\reenc}\big) \;+\; \big(F_e - F_{\reenc}\big)
\;+\; \big(I_e - I_{\reenc}\big) \;+\; \big(R'_e - R'_{\reenc}\big)
&&\text{Eq.~\ref{eq:fourterm}, each arm}\\
&=\;
\big(S_e - S_{\reenc}\big) \;+\; F_e \;+\; I_e
\;+\; \big(R'_e - R'_{\reenc}\big),
&&F_{\reenc} = I_{\reenc} = 0
\end{align*}
so the control collapses to
$d_{\reenc} \approx S_{\reenc} + R'_{\reenc}$, with $S_{\reenc}$ the
re-encode share (measured ${\approx}\,0$ where probed: the control row
of Table~\ref{tab:floor}).
The subtraction cancels only the re-encode component of the data share
(up to a cross term suppressed by $\sqrt{S_{\reenc}}$, bounded by the
control's near-zero reading), so $m$ is henceforth read net of
re-encode; the graph share and the interaction pass through untouched
--- the control has nothing to cancel them with. Disposing of the
surviving terms is what the assumptions below are for.

\paragraph{The reduced two-term account.}
The subtraction form above is exact bookkeeping, but not yet an
instrument: none of its four terms is separately measurable per bin.
Four named assumptions --- each tested by a designated probe in
\S\ref{sec:evidence} --- trade it for a form built from measurable
curves and per-route constants:
(i)~\textbf{small remainder}: $R'_e$ (and its control counterpart)
negligible over the claimed domain;
(ii)~\textbf{negligible projected interaction}:
$|I_e| \ll S_e + F_e$ over the claimed domain;
(iii)~\textbf{graph-relative stationarity}:
$\|(\Delta J^{\!\top}\bar r)^{\perp}\| / \|\bar g_{\src}\|$ is
$\sigma$-constant (numerator and denominator
share the residual's scale), so the graph share $F_e$ is flat in
$\sigma$; write $\Floor_e$ for its constant value;
(iv)~\textbf{data-branch factorization}:
$\|(J^{\!\top}\Delta\bar r)^{\perp}\| \approx a_e\, m(\sigma)$, a
$\sigma$-independent route amplitude times the route-independent
mismatch shape $m(\sigma) = \|\Delta \bar r(\sigma)\|$ of the model's
mean prediction residual across grids (directly measurable,
\S\ref{sec:anatomy}; read net of re-encode per the subtraction above).
Note $m$ is a \emph{norm}: the factorization assumes a
route-independent amplitude curve, not a shared mismatch
\emph{direction} --- and \S\ref{sec:anatomy} finds the directions are
in fact image- and route-specific.
Applied in that order,
\begin{align}
\gap_e(\sigma)
\;&\approx\; \big(S_e - S_{\reenc}\big) \;+\; F_e
&&\text{by (i), (ii)}
\nonumber\\
\;&\approx\;
\underbrace{A_e \cdot \big(m(\sigma)/\|\bar g_{\src}\|\big)^{2}}_{\text{data term}}
\;+\; \underbrace{\Floor_e}_{\text{graph term}},
\qquad A_e = \tfrac{1}{2}\,a_e^2,
&&\text{by (iii), (iv)}
\label{eq:twoterm}
\end{align}
the \emph{two-term account}, stated on the reported form of the
estimand, the re-encoding excess of \S\ref{sec:estimand}, with
$\|\bar g_{\src}(\sigma)\|$ the total gradient norm. The
quadratic power on $m/\|\bar g_{\src}\|$ is fixed by the cosine geometry \emph{locally};
far from the small-mismatch regime the licensed read is the parent form
itself: keeping Eq.~\ref{eq:exact} un-expanded, with the factorization's loading
written as $\kappa_e(\sigma) = c_e \, m(\sigma)/\|\bar g_{\src}(\sigma)\|$ (so
$A_e = c_e^2/2$ in the small-$\kappa$ limit), yields a data term that
\emph{saturates} instead of overshooting when $\kappa$ is comparable to
one. The loading itself (that the orthogonal gradient mismatch is
proportional to the residual mismatch, $\|\delta g^{\perp}\| =
c_e\, m(\sigma)$) is an \emph{empirical} ingredient, not a
derivation; so wherever a fitted functional form is reported, the
un-expanded Eq.~\ref{eq:exact} is labeled the \emph{best empirical
predictor} and Eq.~\ref{eq:twoterm} the \emph{derived
small-perturbation form}; \S\ref{sec:confronted} scores both, held
out.

\paragraph{The graph term has identifiable parts.}
Changing the grid changes (i)~the rotary coordinate system, in which
every band of the position embedding accumulates phase at a different
density across the image, and (ii)~non-positional graph statistics: attention
softmax over a different token count, sequence-length-dependent
normalization, and the sheer capacity of the coarse graph to approximate
the fine graph's computation (the conceptual frame is attention as a
discretization of a function-space operator,
\citealp{calvello2024continuum}, whose statistics need not be
discretization-invariant; the frame supplies no quantitative law for
the dependence). Splitting the graph share
$\Floor_e = \mathrm{RoPE}_e + \mathrm{Resid}_e$ yields one sharp
prediction: the positional share is a \emph{coordinate-system artifact},
so an exact positional-interpolation intervention (evaluating the
demoted grid at fractional rotary coordinates that reproduce the native
relative phase geometry) should remove $\mathrm{RoPE}_e$ wherever the
data branch is inactive ($\sigma \to 1$); and a band-counting sketch
says the residue should grow with absolute grid change (the number of
rotary bands whose per-extent rotation count crosses decorrelation grows
with the extent), attaching an absolute-size governor to
$\mathrm{Resid}_e$. We flag the sketch's epistemic status plainly: like
YaRN's own band thresholds \citep{peng2024yarn}, its constants are
calibrated, not derived; we claim no closed form for $\Floor_e$.

\paragraph{The ledger.}
Explicitly: the exact form (Eq.~\ref{eq:exact}) and the four-term
expansion (Eq.~\ref{eq:fourterm}) are derived unconditionally from
Eqs.~\ref{eq:gapdef} and~\ref{eq:branches}
(Appendix~\ref{app:geometry}); the two-term reduction
(Eq.~\ref{eq:twoterm}) is derived \emph{conditional on the four named
assumptions}, each of which is tested by a designated probe
(\S\ref{sec:evidence}); the
coefficients $A_e$ (or $c_e$), $m(\sigma)$, $\|\bar g_{\src}(\sigma)\|$, $\Floor_e$ are
measured, not derived; and any fitted power or loading is labeled
empirical. The account was
pre-registered with an outcome matrix (including a kill switch: all
endpoint gaps ${\approx}\,0$ would have falsified the graph term) before
the discriminating runs.

\subsection{The spectral account of the data term, from prior work}
\label{sec:null}

The spectral argument reviewed in \S\ref{sec:related} supplies a
second reading of the same gap, and it is worth deriving in the units
the measurements will be stated in. Its strongest available form is
SPD's tolerance-parameterized crossover \citep{xiao2026spd}: under
the premise that each clean-latent frequency $\omega$ is Gaussian,
$x_0^{(\omega)} \sim \mathcal{N}(0, P_\omega)$ (with $P$ the
clean-latent power spectrum), their
Proposition~1 states that for all $\sigma \ge t_\omega$ the
\emph{Bayes-optimal} per-band velocity prediction is within tolerance
$\delta$ of pure noise,
$\mathbb{E}[|v^{*(\omega)} - \epsilon^{(\omega)}|^2] \le \delta$, and
their Proposition~2 places the safe-substitution boundary for a coarse
grid at the $t_\omega$ of its Nyquist band $\omega_e$. Carrying this
to the training gradient --- under the premise that demotion costs
only what the destroyed bands still contribute to the prediction, so
that in the terms of Eq.~\ref{eq:fourterm} the gap is the data share
alone, gated at the Nyquist band's activation time, with no
target-mediated share, no graph share, and no interaction --- the
gradient gap the argument implies is
\begin{equation}
\gap_e^{\mathrm{spec}}(\sigma) \;=\;
S_e(\sigma)\,\mathbf{1}\!\big[\sigma < t_{\omega_e}\big],
\qquad
t_\omega \;=\;
\frac{1}{1 + \sqrt{\delta \,/\, \big(P_\omega\,(1 + P_\omega -
\delta)\big)}},
\qquad
F_e = I_e = 0.
\label{eq:spd}
\end{equation}
This is a prediction-level criterion, strictly stronger than raw input
SNR; the often-quoted equal-power crossover
$\sigma_{\mathrm{eq}}(f) = \sqrt{P(f)}/(1+\sqrt{P(f)})$ is its
$\delta = (1+P_\omega)/2$ slice.

Three structural consequences are fixed by the form of
Eq.~\ref{eq:spd} before any gradient is measured. All route
boundaries move together under the single tolerance $\delta$: the
family is totally ordered by one parameter. In normalized frequency
the demoted grid's Nyquist sits at
$f_{\mathrm{cut}} = \tfrac{1}{2}\,e/e_0$, a function of the downscale
\emph{ratio} alone, so the account cannot distinguish routes of equal
ratio at different absolute sizes and is structurally blind to any
absolute-size effect. And for every route and every $\delta$ the
predicted gap vanishes above a finite boundary: every route becomes
safe at some $\sigma$, and no $\sigma$-independent floor is
expressible. Under the same diagonal-Gaussian premise the account
also extends to the residual level: the posterior mean is a per-band
Wiener shrinkage
(under a Gaussian prior and Gaussian noise the minimum-mean-square-error
estimator is linear with per-band gain $P_\omega/(P_\omega + n_\omega)$,
the classical Wiener filter, \citealp{wiener1949}), so with
no cross-frequency coupling the cross-grid mean-residual mismatch is
confined to the destroyed band and ordered across routes by
destroyed-band energy.

The family is instantiated on our measured latent spectrum in
\S\ref{sec:phenomenon} (Table~\ref{tab:null}); its curve-level read in
gap units necessarily borrows our bridge: the derivation emits residual
units only, so any gap-unit prediction is \emph{the spectral account
transported through} $\|\bar g_{\src}\|$ \emph{and a calibrated route gain}, not a
parameter-free prediction (\S\ref{sec:confronted}). Two properties of
the instantiation are fixed before any gradient is measured: the latents
are high-frequency-quiet, so the family is \emph{compressed} (the
spread between the mildest and harshest route remains below ${\sim}0.13$
in $\sigma$ at every tolerance), and the two routes our measurements
will separate maximally ($1024{\to}896$, ratio $0.875$, and
$896{\to}768$, ratio $0.857$) receive boundaries within $0.01$ of each
other at \emph{every} $\delta$.

\subsection{Finite draws and the debiased estimator}
\label{sec:debias}

The estimand of Eq.~\ref{eq:gapdef} is defined on population mean
gradients; any instrument estimates them from $D$ noise draws. This
finite-draw step is not a neutral nuisance: in principle it biases the
naive estimate in exactly the direction under test, so the correction
belongs to the theory, not to the lab notebook.

\paragraph{Why finite draws manufacture floors.}
The natural estimate of Eq.~\ref{eq:gapdef}, a redraw floor minus a
single demoted-arm cosine, is asymmetric: the floor is a cosine
between \emph{two} finite-draw estimates of the native gradient, while
each demote arm contributes \emph{one} estimate compared against the
native mean. Finite-draw noise attenuates every such cosine below its
population value (each estimated direction carries variance the
population direction lacks), and the attenuation is arm-dependent:
per-draw gradient variance grows as the token count falls, because the
MSE loss averages over fewer elements, so a demoted arm's single
estimate is noisier than the native estimates it is compared to. An
iso-direction null (demotion changing \emph{nothing} but the
variance) therefore still produces a positive gap, and one that
\emph{grows as the target grid shrinks}: the same signature as
``absolute token count sets the floor.'' A token-count-ordered floor
can thus be manufactured entirely by the estimator, and no naive read
can tell the two apart; this is why every verdict in this paper is
stated in debiased units.

\paragraph{The debiased estimator.}
Two corrections, used together. First, \emph{self-floors with
attenuation correction}: every arm runs a second independent draw set
$g_a'$, giving a per-arm self-cosine
$\cos_{\mathrm{self}} = \cos(g_a, g_a')$, so that each arm has the same
two-estimate structure the native floor has. The debiased cosine is
the classical split-half disattenuation correction,
\begin{equation}
\hat c \;=\; \frac{\cos(\bar g_{\src},\, \bar g_{a})}
{\sqrt{\cos_{\mathrm{floor}} \cdot \cos_{\mathrm{self},a}}},
\qquad
\dgap \;=\; 1 - \hat c,
\label{eq:debias}
\end{equation}
which cancels the finite-draw attenuation of numerator and denominator
to first order. Second, \emph{draw-count extrapolation}: the
attenuation of an accumulated-draw cosine decays as $1/D$, so sweeping
nested draw subsets and fitting
$\gap(D) = \gap_\infty + c/D$ reads off the draw-limit gap
$\gap_\infty$ directly. The two corrections are mutually checking: if
the disattenuation removes the finite-draw component, debiased
estimates must come out $D$-flat ($c \approx 0$) under the sweep, a
prediction the instrument verifies (\S\ref{sec:instrument}).

